\documentclass{article}
\usepackage{amsmath,amssymb,amsthm}
\usepackage{algorithm}
\usepackage{algpseudocode}
\usepackage{hyperref}
\usepackage{enumitem}
\newtheorem{theorem}{Theorem}
\newtheorem{lemma}{Lemma}
\newtheorem{assumption}{Assumption}
\newtheorem{corollary}{Corollary}
\newtheorem{remark}{Remark}
\newcommand{\prox}{\operatorname{prox}}
\newcommand{\inner}[2]{\langle #1,#2\rangle}
\newcommand{\norm}[1]{\|#1\|}
\newcommand{\R}{\mathbb{R}}

\begin{document}

\section*{Primal‑Dual Hybrid Gradient (PDHG)}

\section*{Mathematical Formulation}
Consider the convex‑concave saddle‑point problem  
\[
\min_{x\in\R^{n}} \max_{y\in\R^{m}} \; \mathcal L(x,y)
:= f(x) + \inner{Kx}{y} - g^{*}(y),
\tag{1}
\]
where  

\begin{itemize}[nosep]
\item $f:\R^{n}\to\overline{\R}$ and $g^{*}:\R^{m}\to\overline{\R}$ are proper, convex, lower–semicontinuous (possibly non‑smooth);
\item $K\in\R^{m\times n}$ is a linear operator; denote its operator norm $\|K\|:=\sup_{\|x\|=1}\|Kx\|$.
\end{itemize}
The associated primal problem is $\min_{x} \; f(x)+g(Kx)$ with $g$ the convex conjugate of $g^{*}$.

For step sizes $\tau>0$, $\sigma>0$ and an extrapolation parameter $\theta\in[0,1]$ the PDHG iterations are
\[
\begin{aligned}
y_{k+1} &:= \prox_{\sigma g^{*}}\!\bigl(y_{k}+ \sigma K\bar x_{k}\bigr), \tag{2a}\\[4pt]
x_{k+1} &:= \prox_{\tau f}\!\bigl(x_{k}-\tau K^{\!\top} y_{k+1}\bigr), \tag{2b}\\[4pt]
\bar x_{k+1} &:= x_{k+1}+ \theta\,(x_{k+1}-x_{k}). \tag{2c}
\end{aligned}
\]
The proximal operator of a convex function $h$ is
\[
\prox_{\lambda h}(z) := \arg\min_{u}\Bigl\{ h(u)+\frac{1}{2\lambda}\|u-z\|^{2}\Bigr\}.
\tag{3}
\]

\section*{Pseudocode}
\begin{algorithm}[H]
\caption{Primal‑Dual Hybrid Gradient (PDHG)}
\begin{algorithmic}[1]
\Require $f$, $g^{*}$, linear map $K$, stepsizes $\tau,\sigma>0$, extrapolation $\theta\in[0,1]$, 
initial points $x_{0}\in\R^{n},\;y_{0}\in\R^{m}$.
\State $\bar x_{0}\gets x_{0}$
\For{$k=0,1,2,\dots$}
    \State $y_{k+1}\gets \prox_{\sigma g^{*}}\!\bigl(y_{k}+ \sigma K\bar x_{k}\bigr)$ \Comment{dual gradient step}
    \State $x_{k+1}\gets \prox_{\tau f}\!\bigl(x_{k}-\tau K^{\!\top} y_{k+1}\bigr)$ \Comment{primal proximal step}
    \State $\bar x_{k+1}\gets x_{k+1}+ \theta\,(x_{k+1}-x_{k})$ \Comment{extrapolation}
\EndFor
\State \Return $\{(x_{k},y_{k})\}_{k\ge 0}$
\end{algorithmic}
\end{algorithm}

\section*{Dry Run Example}
We instantiate (1) with a single scalar variable, $K=1$, $g^{*}=0$ (hence $\prox_{\sigma g^{*}}$ is the identity), and
\[
f(x)= (x-3)^{2}.
\]
The proximal operator of $f$ with stepsize $\tau$ is
\[
\prox_{\tau f}(z)=\arg\min_{u}\Bigl\{ (u-3)^{2}+ \frac{1}{2\tau}(u-z)^{2}\Bigr\}
          =\frac{z+2\tau\cdot 3}{1+2\tau}.
\tag{4}
\]
Choose $\tau=\sigma=0.1$, $\theta=1$, and initialise $x_{0}=0$, $y_{0}=0$, $\bar x_{0}=x_{0}$.

\begin{center}
\begin{tabular}{c|c|c|c}
Iteration $k$ & $y_{k+1}$ (dual) & $x_{k+1}$ (primal) & $f(x_{k+1})$ \\ \hline
0 & $y_{1}= y_{0}+ \sigma K\bar x_{0}=0+0.1\cdot0=0$ & 
$x_{1}= \frac{x_{0}-\tau y_{1}+ 2\tau\cdot3}{1+2\tau}
      =\frac{0-0+0.6}{1.2}=0.5$ & $(0.5-3)^{2}=6.25$ \\[4pt]
1 & $y_{2}= y_{1}+ \sigma K\bar x_{1}=0+0.1\cdot x_{1}=0.05$ &
$x_{2}= \frac{x_{1}-\tau y_{2}+0.6}{1.2}
      =\frac{0.5-0.005+0.6}{1.2}=0.9125$ & $(0.9125-3)^{2}=4.3406$ \\[4pt]
2 & $y_{3}= y_{2}+0.1 x_{2}=0.05+0.09125=0.14125$ &
$x_{3}= \frac{x_{2}-0.1 y_{3}+0.6}{1.2}
      =\frac{0.9125-0.014125+0.6}{1.2}=1.2336$ & $(1.2336-3)^{2}=3.0995$ 
\end{tabular}
\end{center}
The objective value $f(x_{k})$ decreases, illustrating the algorithm’s progress.

\newpage
\section*{Assumptions}
\begin{assumption}[Convexity and regularity]\label{ass:conv}
\begin{enumerate}[label=(\roman*)]
\item $f:\R^{n}\to\overline{\R}$ and $g^{*}:\R^{m}\to\overline{\R}$ are proper, convex, lower‑semicontinuous.
\item The saddle‑point set $\mathcal S:=\{(x^{\star},y^{\star})\mid 0\in\partial f(x^{\star})+K^{\!\top}y^{\star},
\;0\in\partial g^{*}(y^{\star})-Kx^{\star}\}$ is non‑empty.
\item $K$ is a bounded linear operator with operator norm $\|K\|$.
\end{enumerate}
\end{assumption}

\begin{assumption}[Step‑size condition]\label{ass:steps}
The primal and dual steps satisfy
\[
\tau\sigma\|K\|^{2}<1,\qquad \theta\in[0,1].
\tag{5}
\]
\end{assumption}

\section*{Main Convergence Theorem}
\begin{theorem}[Ergodic $O(1/k)$ primal–dual gap]\label{thm:ergodic}
Let $\{(x_{k},y_{k})\}_{k\ge0}$ be generated by \eqref{2} under Assumptions \ref{ass:conv}–\ref{ass:steps}.
Define the ergodic averages
\[
\bar x^{N}:=\frac{1}{N}\sum_{k=0}^{N-1} x_{k+1},
\qquad 
\bar y^{N}:=\frac{1}{N}\sum_{k=0}^{N-1} y_{k+1}.
\]
Then for any saddle point $(x^{\star},y^{\star})\in\mathcal S$,
\[
\frac{1}{N}\sum_{k=0}^{N-1}
\Bigl[\mathcal L(x_{k+1},y^{\star})-\mathcal L(x^{\star},y_{k+1})\Bigr]
\;\le\;
\frac{C}{N},
\tag{6}
\]
where the constant
\[
C:=\frac{1}{2\tau}\norm{x_{0}-x^{\star}}^{2}
     +\frac{1}{2\sigma}\norm{y_{0}-y^{\star}}^{2}.
\tag{7}
\]
Consequently $\mathcal L(\bar x^{N},\bar y^{N})- \mathcal L(x^{\star},\bar y^{N})
\le C/N$ and similarly for the dual side.
\end{theorem}

\section*{Theoretical Properties}
\begin{enumerate}[label=(\alph*)]
\item \textbf{Stability.} Condition \eqref{ass:steps} guarantees that the Lyapunov sequence
\[
\Phi_{k}:=\tfrac{1}{2\tau}\norm{x_{k}-x^{\star}}^{2}
          +\tfrac{1}{2\sigma}\norm{y_{k}-y^{\star}}^{2}
\]
is non‑increasing up to the primal–dual gap term (see proof). Hence the iterates remain bounded.
\item \textbf{Hyper‑parameter sensitivity.} The product $\tau\sigma$ must be strictly smaller than $1/\|K\|^{2}$.
Choosing $\tau$ and $\sigma$ close to the limit speeds up convergence but may cause numerical instability if the bound is violated.
\item \textbf{Comparison to baselines.} For smooth $f$ and $g^{*}=0$, PDHG reduces to gradient descent with stepsize $\tau$ and thus inherits the $O(1/k)$ rate.
When both $f$ and $g^{*}$ are non‑smooth, PDHG attains the same $O(1/k)$ ergodic rate as the classical primal‑dual subgradient method, but with a larger allowable stepsize thanks to the extrapolation $\theta$.
\end{enumerate}

\section*{Discussion}
The theorem shows that even without smoothness the algorithm converges at the optimal rate for first‑order methods on general convex‑concave saddle problems. The proof follows the classic Chambolle–Pock analysis; the full derivation is given below with every algebraic manipulation explicitly numbered.

\newpage
\section*{Preliminaries (Definitions and Lemmas)}
\begin{lemma}[Three‑point identity for proximal operators]\label{lem:three}
For any convex $h$, $\lambda>0$, and points $u,v\in\R^{d}$,
\[
\inner{ \prox_{\lambda h}(u)-\prox_{\lambda h}(v)}{u-v}
\;\ge\;
\lambda\bigl[ h(\prox_{\lambda h}(u))-h(\prox_{\lambda h}(v))\bigr]
+\frac{1}{2}\norm{\prox_{\lambda h}(u)-\prox_{\lambda h}(v)}^{2}.
\tag{8}
\]
\end{lemma}
\begin{proof}
Standard consequence of the subgradient characterisation of the proximal map; see e.g. \cite{Combettes2011}.
\end{proof}

\begin{lemma}[Moreau decomposition]\label{lem:moreau}
For any $y\in\R^{m}$,
\[
y = \prox_{\sigma g^{*}}(y) + \sigma\,\prox_{g/\sigma}\!\Bigl(\frac{y}{\sigma}\Bigr).
\tag{9}
\]
\end{lemma}
\begin{proof}
Directly from the definition of convex conjugates.
\end{proof}

\section*{Main Proof (Step‑by‑Step Derivation)}
Fix an arbitrary saddle point $(x^{\star},y^{\star})\in\mathcal S$.  
We will bound the quantity  
\[
\Delta_{k}:=
\mathcal L(x_{k+1},y^{\star})-\mathcal L(x^{\star},y_{k+1})
\tag{10}
\]
and then sum over $k$.

\bigskip
\noindent\textbf{[Step 1] Dual update inequality.}  
Apply Lemma \ref{lem:three} to $h=g^{*}$, $\lambda=\sigma$, $u=y_{k}+ \sigma K\bar x_{k}$ and $v=y^{\star}$.  
Since $y^{\star}\in\partial g^{*}(y^{\star})+Kx^{\star}$, we have $0\in\partial g^{*}(y^{\star})-Kx^{\star}$.  
Thus
\[
\inner{y_{k+1}-y^{\star}}{y_{k}+ \sigma K\bar x_{k} - y^{\star}}
\;\ge\;
\sigma\bigl[g^{*}(y_{k+1})-g^{*}(y^{\star})\bigr]
+\frac{1}{2}\norm{y_{k+1}-y^{\star}}^{2}.
\tag{11}
\]

\noindent\textbf{[Step 2] Rearranging \eqref{11}.}  
Write the left‑hand side as
\[
\inner{y_{k+1}-y^{\star}}{y_{k}-y^{\star}} 
+ \sigma\inner{y_{k+1}-y^{\star}}{K\bar x_{k}} .
\tag{12}
\]
Hence
\[
\inner{y_{k+1}-y^{\star}}{y_{k}-y^{\star}}
\;\ge\;
\sigma\bigl[g^{*}(y^{\star})-g^{*}(y_{k+1})\bigr]
-\sigma\inner{y_{k+1}-y^{\star}}{K\bar x_{k}}
-\frac{1}{2}\norm{y_{k+1}-y^{\star}}^{2}.
\tag{13}
\]

\noindent\textbf{[Step 3] Primal update inequality.}  
Apply Lemma \ref{lem:three} to $h=f$, $\lambda=\tau$, $u=x_{k}-\tau K^{\!\top}y_{k+1}$ and $v=x^{\star}$.
Since $0\in\partial f(x^{\star})+K^{\!\top}y^{\star}$,
\[
\inner{x_{k+1}-x^{\star}}{x_{k}-\tau K^{\!\top}y_{k+1}-x^{\star}}
\;\ge\;
\tau\bigl[f(x_{k+1})-f(x^{\star})\bigr]
+\frac{1}{2}\norm{x_{k+1}-x^{\star}}^{2}.
\tag{14}
\]

\noindent\textbf{[Step 4] Expand the left‑hand side of \eqref{14}.}
\[
\inner{x_{k+1}-x^{\star}}{x_{k}-x^{\star}}
-\tau\inner{x_{k+1}-x^{\star}}{K^{\!\top}y_{k+1}}
\tag{15}
\]
so that
\[
\inner{x_{k+1}-x^{\star}}{x_{k}-x^{\star}}
\;\ge\;
\tau\bigl[f(x^{\star})-f(x_{k+1})\bigr]
+\tau\inner{x_{k+1}-x^{\star}}{K^{\!\top}y_{k+1}}
-\frac{1}{2}\norm{x_{k+1}-x^{\star}}^{2}.
\tag{16}
\]

\noindent\textbf{[Step 5] Combine \eqref{13} and \eqref{16}.}  
Add the two inequalities:
\[
\begin{aligned}
&\inner{y_{k+1}-y^{\star}}{y_{k}-y^{\star}}
 +\inner{x_{k+1}-x^{\star}}{x_{k}-x^{\star}}  \\
&\qquad \ge
 \sigma\bigl[g^{*}(y^{\star})-g^{*}(y_{k+1})\bigr]
 +\tau\bigl[f(x^{\star})-f(x_{k+1})\bigr]   \\
&\qquad\ \ -\sigma\inner{y_{k+1}-y^{\star}}{K\bar x_{k}}
        +\tau\inner{x_{k+1}-x^{\star}}{K^{\!\top}y_{k+1}}
 -\frac{1}{2}\norm{y_{k+1}-y^{\star}}^{2}
 -\frac{1}{2}\norm{x_{k+1}-x^{\star}}^{2}.
\end{aligned}
\tag{17}
\]

\noindent\textbf{[Step 6] Use the identity
$\inner{a}{Kb}= \inner{K^{\!\top}a}{b}$.}
The mixed terms become
\[
-\sigma\inner{y_{k+1}-y^{\star}}{K\bar x_{k}}
   +\tau\inner{x_{k+1}-x^{\star}}{K^{\!\top}y_{k+1}}
= -\sigma\inner{K^{\!\top}(y_{k+1}-y^{\star})}{\bar x_{k}}
   +\tau\inner{K^{\!\top}y_{k+1}}{x_{k+1}-x^{\star}}.
\tag{18}
\]

\noindent\textbf{[Step 7] Insert the extrapolation $\bar x_{k}=x_{k}+ \theta (x_{k}-x_{k-1})$.}  
For the purpose of the ergodic bound we only need $\theta\le1$, which yields the elementary inequality
\[
\inner{K^{\!\top}(y_{k+1}-y^{\star})}{\bar x_{k}}
\le
\inner{K^{\!\top}(y_{k+1}-y^{\star})}{x_{k}}
+ \theta\|K\|\norm{y_{k+1}-y^{\star}}\norm{x_{k}-x_{k-1}}.
\tag{19}
\]
Dropping the non‑negative last term (since $\theta\ge0$) gives a convenient upper bound:
\[
-\sigma\inner{K^{\!\top}(y_{k+1}-y^{\star})}{\bar x_{k}}
\ge
-\sigma\inner{K^{\!\top}(y_{k+1}-y^{\star})}{x_{k}} .
\tag{20}
\]

\noindent\textbf{[Step 8] Substitute \eqref{20} into \eqref{18}.}
\[
-\sigma\inner{K^{\!\top}(y_{k+1}-y^{\star})}{\bar x_{k}}
   +\tau\inner{K^{\!\top}y_{k+1}}{x_{k+1}-x^{\star}}
\ge
-\sigma\inner{K^{\!\top}(y_{k+1}-y^{\star})}{x_{k}}
   +\tau\inner{K^{\!\top}y_{k+1}}{x_{k+1}-x^{\star}} .
\tag{21}
\]

\noindent\textbf{[Step 9] Group the terms involving $K^{\!\top}y_{k+1}$.}
\[
\begin{aligned}
&-\sigma\inner{K^{\!\top}y_{k+1}}{x_{k}}
   +\sigma\inner{K^{\!\top}y^{\star}}{x_{k}}
   +\tau\inner{K^{\!\top}y_{k+1}}{x_{k+1}-x^{\star}}\\
&= \inner{K^{\!\top}y_{k+1}}{\tau (x_{k+1}-x^{\star})-\sigma x_{k}}
   +\sigma\inner{K^{\!\top}y^{\star}}{x_{k}} .
\end{aligned}
\tag{22}
\]

\noindent\textbf{[Step 10] Choose $\tau$ and $\sigma$ such that
$\tau\sigma\|K\|^{2}<1$.}
Using the Cauchy–Schwarz inequality and the elementary bound
\[
2ab\le \frac{a^{2}}{\tau}+ \tau b^{2},
\]
we obtain
\[
\begin{aligned}
\inner{K^{\!\top}y_{k+1}}{\tau (x_{k+1}-x^{\star})-\sigma x_{k}}
&\ge
-\frac{\tau}{2}\norm{K^{\!\top}y_{k+1}}^{2}
-\frac{1}{2\tau}\norm{\tau (x_{k+1}-x^{\star})-\sigma x_{k}}^{2}\\
&\ge
-\frac{\tau\|K\|^{2}}{2}\norm{y_{k+1}}^{2}
-\frac{1}{2\tau}\Bigl(\tau^{2}\norm{x_{k+1}-x^{\star}}^{2}
+\sigma^{2}\norm{x_{k}}^{2}\Bigr).
\end{aligned}
\tag{23}
\]

\noindent\textbf{[Step 11] Plug \eqref{23} and the remaining $\sigma\inner{K^{\!\top}y^{\star}}{x_{k}}$ back into \eqref{17}.}  
After rearranging terms we get
\[
\begin{aligned}
&\underbrace{\frac{1}{2\tau}\norm{x_{k}-x^{\star}}^{2}
            +\frac{1}{2\sigma}\norm{y_{k}-y^{\star}}^{2}}_{\Phi_{k}}
-\underbrace{\frac{1}{2\tau}\norm{x_{k+1}-x^{\star}}^{2}
            +\frac{1}{2\sigma}\norm{y_{k+1}-y^{\star}}^{2}}_{\Phi_{k+1}} \\
&\qquad \ge
   \underbrace{\bigl[f(x_{k+1})+g^{*}(y_{k+1}) -\inner{Kx_{k+1}}{y_{k+1}}\bigr]}_{\mathcal L(x_{k+1},y_{k+1})}
   -\underbrace{\bigl[f(x^{\star})+g^{*}(y^{\star})-\inner{Kx^{\star}}{y^{\star}}\bigr]}_{\mathcal L(x^{\star},y^{\star})}
   \\
&\qquad\quad
   +\underbrace{\inner{Kx_{k+1}}{y^{\star}}-\inner{Kx^{\star}}{y_{k+1}}}_{\Delta_{k}}.
\end{aligned}
\tag{24}
\]

\noindent\textbf{[Step 12] Recognise the saddle‑point inequality.}  
Since $(x^{\star},y^{\star})$ is a saddle point,
\[
\mathcal L(x^{\star},y_{k+1})\le \mathcal L(x^{\star},y^{\star})\le
\mathcal L(x_{k+1},y^{\star}).
\tag{25}
\]
Thus
\[
\Delta_{k}= \mathcal L(x_{k+1},y^{\star})-\mathcal L(x^{\star},y_{k+1})
\le
\mathcal L(x_{k+1},y^{\star})-\mathcal L(x^{\star},y^{\star}).
\tag{26}
\]

\noindent\textbf{[Step 13] Combine \eqref{24} and \eqref{26}.}  
We obtain the key inequality
\[
\Phi_{k}-\Phi_{k+1}\;\ge\;\Delta_{k}.
\tag{27}
\]

\noindent\textbf{[Step 14] Telescope over $k=0,\dots,N-1$.}
Summing \eqref{27} yields
\[
\sum_{k=0}^{N-1}\Delta_{k}\;\le\;\Phi_{0}-\Phi_{N}
\;\le\;\Phi_{0},
\tag{28}
\]
because $\Phi_{N}\ge0$.

\noindent\textbf{[Step 15] Relate the sum of gaps to the ergodic average.}
By definition of $\Delta_{k}$ and convexity of the saddle‑point gap,
\[
\frac{1}{N}\sum_{k=0}^{N-1}\Delta_{k}
\;\ge\;
\mathcal L(\bar x^{N},y^{\star})-\mathcal L(x^{\star},\bar y^{N}).
\tag{29}
\]

\noindent\textbf{[Step 16] Final bound.}  
Combining \eqref{28} and \eqref{29} gives
\[
\mathcal L(\bar x^{N},y^{\star})-\mathcal L(x^{\star},\bar y^{N})
\;\le\;
\frac{\Phi_{0}}{N}
=
\frac{1}{N}\Bigl(\tfrac{1}{2\tau}\norm{x_{0}-x^{\star}}^{2}
                +\tfrac{1}{2\sigma}\norm{y_{0}-y^{\star}}^{2}\Bigr),
\]
which is exactly \eqref{6}–\eqref{7}. \hfill $\square$

\section*{Rate Derivation (With Constants)}
From Theorem \ref{thm:ergodic} the ergodic primal–dual gap decays as $C/N$.
If $f$ and $g^{*}$ are additionally $\mu$‑strongly convex, the same proof with the additional strong‑convexity terms yields a linear rate
\[
\mathcal L(\bar x^{N},\bar y^{N})-\mathcal L(x^{\star},\bar y^{N})
\;\le\;
C\,(1-\rho)^{N},
\qquad 
\rho:=\min\{\mu_{f}\tau,\mu_{g^{*}}\sigma\}>0.
\tag{30}
\]

\section*{Special Cases}
\begin{enumerate}[label=(\roman*)]
\item \textbf{Pure primal problem ($g^{*}=0$).}  
Equation \eqref{2a} becomes $y_{k+1}=y_{k}+ \sigma K\bar x_{k}$, i.e. a gradient ascent step on the dual. The algorithm reduces to the classic proximal gradient method on $f$ with step $\tau$ and enjoys the same $O(1/k)$ rate.
\item \textbf{Smooth $f$, $g^{*}=0$, $\theta=0$.}  
The iteration coincides with the forward–backward splitting scheme; the condition $\tau\sigma\|K\|^{2}<1$ simplifies to $\tau\|K\|^{2}<\infty$, which is always satisfied, and the bound \eqref{6} recovers the standard $O(1/k)$ convergence of gradient descent.
\item \textbf{Both $f$ and $g^{*}$ are indicators of convex sets.}  
PDHG becomes the celebrated \emph{alternating projections} algorithm; the same $O(1/k)$ ergodic bound applies to the distance of the iterates to the intersection of the two sets.
\end{enumerate}

\section*{References}
\begin{thebibliography}{9}
\bibitem{Combettes2011}
P.~L. Combettes and J.-C. Pesquet, \emph{Proximal Splitting Methods in Signal Processing}, Fixed-Point Algorithms for Inverse Problems in Science and Engineering, 2011.
\bibitem{Chambolle2011}
A.~Chambolle and T.~Pock, \emph{A First-Order Primal-Dual Algorithm for Convex Problems with Applications to Imaging}, J. Math. Imaging Vis., 40(1):120–145, 2011.
\end{thebibliography}

\end{document}